\section{Discussion}
\label{sec:discussion}

\subsection{Paper Configuration}
\label{sec:discussion:config}

We select \texttt{avmp\_dynamic\_b128} as the default configuration rather than the \texttt{b256} variant. While \texttt{b256} records marginally fewer OOM events (509.3 versus 510.0 cross-workload), this 0.7 OOM gap falls completely within the per-cell standard deviation of 0.8 to 3.0. Furthermore, the \texttt{b128} configuration migrates 313 MiB per cell on average compared to 352 MiB for \texttt{b256}, representing a 12.5\% reduction in migration churn. We choose \texttt{b128} as the conservative point in the near-optimal region.

\subsection{Future Work}
\label{sec:discussion:future}

Future research directions extend beyond resolving the immediate architectural limitations. We plan to explore workload-prediction heuristics for proactive capacity rebalancing before allocation exceptions occur. Triton kernel integration and a third heterogeneous pool (e.g., dedicated KV-prefix-cache) are planned extensions.

\subsection{Production Hypothesis}
\label{sec:discussion:hypothesis}

We frame the production impact of AVMP as a hypothesis to be tested, not a measured outcome. A Jamba 1.5 Mini deployment requires approximately 24 GiB VRAM per H100 80GB instance for model weights and activations, leaving approximately 56 GiB for KV cache and SSM state combined. Under static dual-pool partitioning at the SGLang default \texttt{mamba\_full\_memory\_ratio} of 0.9, our cross-workload data records 1221 OOM events versus 552 OOMs at the \texttt{mr=0.5} ratio, indicating that default static configurations underutilize one pool while the other saturates. Our dynamic rebalancing reduces OOM by an additional 7.6\% (510 versus 552 cross-workload) over the best static configuration. If production traffic exhibits similar asymmetric pool slack to our synthetic workloads, the OOM reduction suggests possible capacity gains. Validating this hypothesis requires production trace replay (e.g., ShareGPT, Alpaca) and a \texttt{cuMemMap}-backed implementation to remove the current 2$\times$ VRAM overhead. OOM count reduction does not directly translate to concurrent sequence capacity, which depends on scheduler admission policy, prefill/decode split, latency SLO, output length distribution, and model kernel behavior outside our synthetic evaluation scope. The 13.3$\times$ goodput improvement on \texttt{uniform\_short} suggests that throughput gains could exceed the OOM reduction in short-prompt-dominated traffic, though this requires production trace validation. We do not project absolute dollar savings here because real-world cost depends on instance pricing, utilization rates, and workload mix, all outside our synthetic evaluation scope.
